\subsubsection{Gauduchon polar trace comparisons}
\label{subsec:gauduchon-polar-comparison}

Gauduchon trace gives a uniform normalization for polar curvature on different modifications. The resulting comparison separates analytic models without a bound on top-degree mass. Its complete coordinate space isolates the further realization question.

Let $X$ be connected and compact, of dimension $n\geq1$, and let $\omega$ be Hermitian. Fix a set $\mathscr R$ of smooth modifications $\pi:Y_\pi\to X$ containing the identity and common log resolutions for pairs of analytic singularity types. Choose a Gauduchon form $\eta_\pi$ on each $Y_\pi$. For genuine $u\in\PSH(X,\omega)$, set
\[
 \mu_\pi(u)=\mathbf1_{\{u\circ\pi=-\infty\}}
       \bigl(\pi^*\omega+\ddc(u\circ\pi)\bigr)
          \wedge\eta_\pi^{n-1},
\]
with coefficientwise current restriction, and define
\begin{equation}\label{eq:gauduchon-polar-comparison}
 \begin{gathered}
 A_\pi=\int_{Y_\pi}\pi^*\omega\wedge\eta_\pi^{n-1},\\
 \delta_\omega^{\mathrm{pol}}(u,v)
       =\sup_{\pi\in\mathscr R}
          \frac{\|\mu_\pi(u)-\mu_\pi(v)\|_{\mathrm{TV}}}{A_\pi}.
 \end{gathered}
\end{equation}
The Gauduchon forms are fixed data. Positive scalar rescaling of any one of them leaves its coordinate comparison unchanged.

\Needspace{16\baselineskip}
\begin{proposition}\label{prop:gauduchon-polar-comparison}\index{Model potential!Gauduchon polar trace}\index{Gauduchon metric!polar trace comparison}
Every $A_\pi$ is positive. The comparison in \eqref{eq:gauduchon-polar-comparison} is a pseudometric on $\PSH(X,\omega)$, bounded by two, and restricts to a metric on normalized ordinary models with analytic singularities. On all ordinary curve models,
\[
 \delta_\omega^{\mathrm{pol}}(p,q)
       =\frac{\|\mu_{\operatorname{id}_X}(p)-\mu_{\operatorname{id}_X}(q)\|_{\mathrm{TV}}}{\int_X\omega},
\]
so it induces the complete curve topology with analytic closure exactly the ideal models.

For real smooth $f$, put $\omega'=\omega+\ddc f$. Using the same Gauduchon forms and the same defining formula for $\omega'$, its denominators remain positive and
\[
 \delta_{\omega'}^{\mathrm{pol}}(u-f,v-f)
       =\delta_\omega^{\mathrm{pol}}(u,v)
       \qquad(u,v\in\PSH(X,\omega)).
\]
\end{proposition}
\begin{proof}
The form $\pi^*\omega$ is semipositive and is positive on the nonempty isomorphism locus, so $A_\pi>0$. Pullback is genuine and admissible by \cref{prop:psh-descent-modification}. Since $\ddc\eta_\pi^{n-1}=0$, distributional Stokes gives
\[
 \int_{Y_\pi}\bigl(\pi^*\omega+\ddc(u\circ\pi)\bigr)
                 \wedge\eta_\pi^{n-1}=A_\pi.
\]
Thus $\mu_\pi(u)$ is positive and has mass at most $A_\pi$. The normalized variation distance of two such measures is at most two. Its triangle inequality and symmetry persist under the supremum, giving the pseudometric assertion.

For analytic ordinary models $p,q$, choose a common log resolution in $\mathscr R$. If $\delta_\omega^{\mathrm{pol}}(p,q)=0$, their measures in this coordinate agree. If $E_i$ are its resolved prime components and $a_i,b_i$ their coefficients, the divisorial-restriction argument of \cref{lem:nonclosed-divisorial-removal} gives
\[
 \mu_\pi(p)=\sum_i a_i[E_i]\wedge\eta_\pi^{n-1},\qquad
 \mu_\pi(q)=\sum_i b_i[E_i]\wedge\eta_\pi^{n-1}.
\]
A smooth local correction makes the Hessian closed and positive without changing its coefficient mass on a divisor. Distinct components have no overlapping trace mass on their intersections. Each component has strictly positive trace mass. Equality of these measures therefore implies $a_i=b_i$ for every $i$. The finite divisorial criterion of \cref{lem:analytic-ideal-dominance}, applied in both directions, gives $p\sim q$. Ordinary projection depends only on this bounded equivalence class, so normalized ordinary models in that class are equal. This proves separation.

On a curve every smooth modification is an isomorphism, and the trace exponent is zero. Hence $A_\pi=\int_X\omega$, and pullback preserves polar total variation. The displayed curve formula and its assertions follow from \cref{cor:curve-model-completeness,prop:curve-analytic-distance-defect}.

Finally $u-f$ has the same pulled-back curvature current and pole set in background $\omega'$ as $u$ in background $\omega$. Stokes also gives $A_{\pi,\omega'}=A_{\pi,\omega}$. The gauge formula follows. It applies to the genuine potentials before ordinary-model projection.
\end{proof}

The positive measure coordinates have a complete ambient space. For each $\pi$, let $\mathscr B_\pi$ consist of the positive finite Borel measures on $Y_\pi$ with mass at most $A_\pi$.

\Needspace{10\baselineskip}
\begin{proposition}\label{prop:polar-trace-coordinate-completeness}\index{Metric completeness!polar trace coordinates}
The space $\prod_{\pi\in\mathscr R}\mathscr B_\pi$ is complete for the metric
\[
 \widehat\delta((\lambda_\pi),(\nu_\pi))
       =\sup_{\pi\in\mathscr R}
          \frac{\|\lambda_\pi-\nu_\pi\|_{\mathrm{TV}}}{A_\pi}.
\]
The map $u\mapsto(\mu_\pi(u))_\pi$ takes values in this space and preserves the comparison $\delta_\omega^{\mathrm{pol}}$.
\end{proposition}
\begin{proof}
The metric is finite, bounded by two, and separates coordinates. A Cauchy sequence is Cauchy in total variation in each coordinate. Completeness of finite signed measures gives its limit there; positivity and the mass bound pass to that limit. Given $\varepsilon>0$, choose $N$ so that for all $j,k\geq N$ every normalized coordinate distance is at most $\varepsilon$. Letting $k$ tend to infinity separately in each coordinate retains that inequality. Taking the supremum proves convergence in $\widehat\delta$. The final assertion follows from the definition and the mass bound already proved.
\end{proof}

Bounded equivalence also preserves these coordinates, by \cref{prop:bounded-type-polar-invariance}. This allows the raw gauge identity to pass through the ordinary projection needed for model normalization.

\Needspace{12\baselineskip}
\begin{corollary}\label{cor:polar-model-gauge-transport}\index{Model potential!polar gauge transport}
The comparison $\delta_\omega^{\mathrm{pol}}$ descends to ordinary bounded-difference classes. For real smooth $f$ and $\omega'=\omega+\ddc f$, the map
\[
 G_f(p)=P_{\omega'}[p-f]
\]
is a bijection between the normalized ordinary models in the two backgrounds, with inverse $q\mapsto P_\omega[q+f]$. It preserves $\delta^{\mathrm{pol}}$, analytic models and ideal models, using the same Gauduchon forms in the two backgrounds.
\end{corollary}
\begin{proof}
Boundedly equivalent potentials have boundedly equivalent pullbacks on every modification, with the same pole set. \Cref{prop:bounded-type-polar-invariance} therefore gives equality of their polar currents and measures $\mu_\pi$ in every coordinate. This proves descent to ordinary classes.

If $p=P_\omega[p]$, the ordinary version of \cref{prop:i-envelope-gauge} gives
\[
 p+\min_Xf\leq G_f(p)+f\leq p+\max_Xf.
\]
Thus $G_f(p)\sim p-f$. Since ordinary envelopes depend only on bounded equivalence classes, $P_{\omega'}[G_f(p)]=P_{\omega'}[p-f]=G_f(p)$. Also $P_\omega[G_f(p)+f]=p$. Applying the same argument to any model in background $\omega'$ proves surjectivity and the inverse formula. Polar invariance and the raw gauge identity in \cref{prop:gauduchon-polar-comparison} now give
\[
 \delta_{\omega'}^{\mathrm{pol}}(G_f(p),G_f(q))
       =\delta_\omega^{\mathrm{pol}}(p,q).
\]

Smooth subtraction and bounded equivalence preserve analytic singularities. If $p=P_\omega[p]_{\mathcal I}$, the ideal gauge estimate gives $P_{\omega'}[p-f]_{\mathcal I}\sim p-f\sim G_f(p)$. The ideal envelope is an ordinary fixed point by \cref{cor:ordinary-preserves-i}. Two ordinary fixed points with the same bounded equivalence class coincide. Hence it equals $G_f(p)$, proving preservation of ideal models; the inverse gives the converse.
\end{proof}

The scalar traces retain the whole polar currents. Their closedness follows from the local cutoff argument of \cref{prop:polar-curvature-closedness}, including nonclosed pole sets.

\Needspace{10\baselineskip}
\begin{corollary}\label{cor:closed-pole-trace-recovery}\index{Polar curvature!recovery from trace}
Let $u,v\in\PSH(X,\omega)$. For every $\pi\in\mathscr R$, their coefficientwise polar curvature currents are closed and agree if and only if $\mu_\pi(u)=\mu_\pi(v)$. Consequently, $\delta_\omega^{\mathrm{pol}}(u,v)=0$ exactly when all these polar currents agree.
\end{corollary}
\begin{proof}
Pullback is genuine and admissible by \cref{prop:psh-descent-modification}. The two polar currents are closed and positive by \cref{prop:polar-curvature-closedness}, applied with background $\pi^*\omega$. Both are carried by the union of the two Borel pole sets, which has Lebesgue measure zero. Equality of their $\eta_\pi$-traces gives equality of the currents by \cref{lem:closed-null-current-trace}; the converse is immediate. Positivity of every $A_\pi$ proves the last assertion.
\end{proof}

For a comparison-Cauchy sequence, convergence of the traces prevents limiting polar mass from leaving the union of its pole sets. Trace recovery then gives compatible limits of the full currents.

\Needspace{12\baselineskip}
\begin{proposition}\label{prop:closed-pole-compatible-limits}\index{Polar curvature!compatible current limits}
Let $(p_j)$ be normalized ordinary $\omega$-models, Cauchy for $\delta_\omega^{\mathrm{pol}}$. For every $\pi\in\mathscr R$, their polar currents converge weakly to a unique closed positive current $R_\pi$, with
\[
 R_\pi\wedge\eta_\pi^{n-1}
       =\lim_{j\to\infty}\mu_\pi(p_j)
       \quad\text{in total variation}.
\]
If $g:Y_{\pi'}\to Y_\pi$ is a modification with $\pi'=\pi\circ g$ and $\pi,\pi'\in\mathscr R$, then $g_*R_{\pi'}=R_\pi$. There is a genuine $h\in\PSH(X,\omega)$ such that every $R_\pi$ is carried coefficientwise by $\{h\circ\pi=-\infty\}$.
\end{proposition}
\begin{proof}
Fix $\pi$. By \cref{prop:polar-trace-coordinate-completeness}, the measures $\mu_\pi(p_j)$ converge in total variation to a positive measure $\lambda_\pi$. Put
\[
 B_\pi=\bigcup_{j\geq1}\{p_j\circ\pi=-\infty\}.
\]
This is a Borel set of Lebesgue measure zero. Every coordinate measure vanishes on its complement, so $\lambda_\pi(Y_\pi\setminus B_\pi)=0$. The polar currents are closed and positive by \cref{cor:closed-pole-trace-recovery}; their $\eta_\pi$-trace masses are at most $A_\pi$. Compactness of positive currents with bounded mass gives a weakly convergent subsequence. Its limit $R$ is closed and positive, and $R\wedge\eta_\pi^{n-1}=\lambda_\pi$. Positivity bounds all coefficient measures by the Hermitian trace, so $R$ is carried by $B_\pi$. Any two subsequential limits have this same null carrier and trace. They agree by \cref{lem:closed-null-current-trace}. Every subsequence admits a further convergent subsequence, hence the whole current sequence converges weakly to this unique $R_\pi$.

Apply \cref{prop:polar-curvature-pushforward} to $g$ with background $\pi^*\omega$ and potential $p_j\circ\pi$. It gives $g_*\operatorname{Pol}_{\pi'}(p_j)=\operatorname{Pol}_\pi(p_j)$. Proper pushforward is weakly continuous, since testing downstairs pulls back a smooth test form upstairs. Passing to the limits proves compatibility. The limits are constructed separately in each coordinate; no diagonal subsequence over $\mathscr R$ is required.

To obtain one pole carrier, put $h_m=\sum_{j=1}^m2^{-j}p_j$. These functions decrease and are $\omega$-plurisubharmonic, because the coefficients have sum at most one and $\omega>0$. Normalized compactness in \cref{thm:normalized-compactness} bounds their $L^1$ norms uniformly. Thus \cref{cor:decreasing-limit} gives a genuine limit $h\in\PSH(X,\omega)$. Each pole set of $p_j$ is contained in the pole set of $h$, since all summands are nonpositive. Hence $B_\pi\subset\{h\circ\pi=-\infty\}$, giving the last assertion.
\end{proof}

The common pole carrier does not identify $R_\pi$ with the polar curvature of $h$. Realization of all these limits by one ordinary model is a further problem.

The choice of Gauduchon forms also affects convergence. On a surface, a form close to a pulled-back K\"ahler form assigns little trace mass to exceptional curves. Choosing these forms independently can suppress a sequence of moving point poles of fixed strength. We compare with the upper-volume comparison $d_\omega$ and the positive lower-volume sectors $\mathcal M_\delta(\omega)$ of \cref{def:model-volume-comparison}.

\Needspace{12\baselineskip}
\begin{proposition}\label{prop:gauduchon-choice-topology-obstruction}\index{Gauduchon metric!choice and model topology}\index{Model potential!moving point poles}
Let $X=\mathbb P^2$ and let $\omega$ be the Fubini--Study form normalized by $\int_X\omega^2=1$. For every set $\mathscr R$ as above and every rational $a\in(0,1)$, there are fixed Gauduchon forms $\eta_\pi$ and analytic ideal models $p_j$ with pairwise distinct single-point pole sets, each of logarithmic coefficient $a$, such that
\[
 \delta_\omega^{\mathrm{pol}}(p_j,0)\leq\frac1j,
 \qquad d_\omega(p_j,0)=a^2.
\]
All $p_j$ and $0$ belong to $\mathcal M_{(1-a^2)/2}(\omega)$. Thus these allowed choices of the forms do not induce the positive-volume topology of \cref{thm:positive-model-complete-metric}.
\end{proposition}
\begin{proof}
At $x=[1:0:0]$, put
\[
 u_x([Z_0:Z_1:Z_2])
   =a\log\frac{|Z_1|^2+|Z_2|^2}
                    {|Z_0|^2+|Z_1|^2+|Z_2|^2}.
\]
The homogeneous quotient defines a global nonpositive potential. Its curvature is positive, since the logarithm of the numerator has positive Hessian in a local projective trivialization and the remaining background is $(1-a)\omega$. Its only pole is $x$, with local expression $a\log|z|^2+O(1)$. By \cref{prop:analytic-envelope-type}, $p_x=P_\omega[u_x]$ is an analytic ideal model ordinarily equivalent to $u_x$. Unitary projective automorphisms preserve $\omega$ and commute with its zero-obstacle envelopes. Transport this construction to pairwise distinct points $x_j$ and write the resulting models as $p_j$.

Fix $\pi:Y_\pi\to X$ and an initial Gauduchon form $\gamma_\pi$. The closed polar current of $p_j\circ\pi$ is supported on the analytic fiber over $x_j$. The divisorial restriction in \cref{lem:nonclosed-divisorial-removal} and the support theorem, \cite[Chapter III, Corollaries 2.11 and 2.14, pp.~141--143]{DemCADG}, give
\[
 \operatorname{Pol}_\pi(p_j)
    =\sum_{E\mapsto x_j}b_{E,j}[E],
 \qquad b_{E,j}=a\operatorname{ord}_E
                         (\mathfrak m_{x_j}\mathcal O_{Y_\pi}).
\]
Indeed, \cref{cor:closed-pole-trace-recovery} supplies closedness, and after removing the curve components the remaining closed order-zero current is supported at finitely many points and vanishes. The sum is finite. Every such curve is contracted by $\pi$, so the restriction of $\pi^*\omega$ to it is zero. Put
\[
 C_{\pi,j}=\sum_{E\mapsto x_j}b_{E,j}\int_E\gamma_\pi,
 \qquad J_\pi=\{j:C_{\pi,j}>0\}.
\]
A surface modification has finitely many exceptional curves and hence finitely many images of such curves. Since the points $x_j$ are distinct, $J_\pi$ is finite. Choose $\varepsilon_\pi>0$ with $\varepsilon_\pi C_{\pi,j}\leq1/j$ for every $j\in J_\pi$; if $J_\pi$ is empty, take $\varepsilon_\pi=1$. Now fix
\[
 \eta_\pi=\pi^*\omega+\varepsilon_\pi\gamma_\pi.
\]
\Needspace{6\baselineskip}
This is positive and Gauduchon: on a surface its Gauduchon equation is $\ddc\eta_\pi=0$. Moreover
\[
 A_\pi=1+\varepsilon_\pi
                  \int_{Y_\pi}\pi^*\omega\wedge\gamma_\pi\geq1,
 \qquad \|\mu_\pi(p_j)\|_{\mathrm{TV}}
                  =\varepsilon_\pi C_{\pi,j}\leq\frac1j.
\]
These choices are fixed for the whole sequence, and the estimate holds for every modification. Since $\mu_\pi(0)=0$, taking the supremum gives the asserted convergence.

To compute singularity volumes, let $\rho:Y\to X$ be the blowup at $x$, with exceptional curve $E$, and let $F:Y\to\mathbb P^1$ resolve the pencil $[Z_1:Z_2]$. Write $\omega_1$ for the normalized Fubini--Study form on $\mathbb P^1$. The current identity is
\[
 \rho^*\omega+\ddc(u_x\circ\rho)=a[E]+H_a,
 \qquad H_a=(1-a)\rho^*\omega+aF^*\omega_1.
\]
For every $v\sim u_x$, divisorial removal gives $H_a+\ddc w\geq0$ with $w=v\circ\rho-u_x\circ\rho$ bounded after extension across $E$. The bounded product of \cref{prop:bounded-products} assigns no mass to $E$. Change of variables off $E$ therefore identifies the non-pluripolar mass of $v$ with $\int_Y(H_a+\ddc w)^2$. Since $H_a$ is closed, bounded Bedford--Taylor integration by parts gives
\[
 \int_Y(H_a+\ddc w)^2=\int_YH_a^2=1-a^2.
\]
For the last equality, the residual class is $H-a[E]$, where $H=\rho^*[\omega]$, with $H^2=1$, $H\cdot E=0$ and $E^2=-1$. Hence both singularity volumes of $p_x$, and by unitary transport of every $p_j$, are $1-a^2$. The zero model has both volumes equal to one. Since $p_j\leq0$, \cref{def:model-volume-comparison} gives $d_\omega(p_j,0)=1-(1-a^2)=a^2$. This also proves the common positive lower-volume bound and the topological incompatibility.
\end{proof}

General model separation, realization of coordinate limits and exact analytic closure remain unresolved. The preceding obstruction shows that compatibility with the positive-volume topology requires more than an arbitrary independent choice of Gauduchon forms. This normalization does not treat backgrounds of zero curve area.
